
\begin{equation}\tag{1.13.2}
  x^{r+1}=0.
\end{equation}

\textbf{Remarque 1.14.} En utilisant la notion de «schéma relatif» (cf [3]), on peut encore donner un sens à 1.1 en supposant que \(S\) est un topos localement annelé. Si on suppose que l'objet final du topos est quasi-compact, la démonstration donnée dans le cas des schémas se transpose sans difficultés, et le théorème reste donc valable.

\section*{2. Calcul du \(K^{\bullet}\) d'un fibré projectif : cas des complexes parfaits}


Le présent paragraphe ne sera pas utilisé dans la suite du séminaire.

La notation \(K^{\bullet}(X)\) désignera dans ce paragraphe le groupe de Grothendieck des complexes parfaits sur \(X\) (IV 2.2). Nous nous proposons de prouver le théorème suivant, analogue de 1.1 :

\textbf{Théorème 2.1.} \uline{Soient} \(S\) \uline{un schéma quasi-compact}, \(\underline{E}\) \uline{un} \(\underline{O}_S\)\uline{-Module localement libre de rang} \(r+1\), \(X=\mathbf P(\underline{E})\) \uline{le fibré projectif associé} ; \uline{soit} \(\xi\) \uline{la classe} \uline{dans} \(K^{\bullet}(X)\) \uline{du complexe égal à} \(\underline{O}_X(1)\) \uline{en degré nul}, \uline{et nul pour les autres degrés}. \uline{Alors} \(K^{\bullet}(X)\), \uline{muni de sa structure naturelle de} \(K^{\bullet}(S)\)-\uline{algèbre est} \uline{un} \(K^{\bullet}(S)\)-\uline{module libre admettant pour base} \(1,\xi,\ldots,\xi^r\).

2.2. Fixons les notations : \(\underline{B}\) désignera l'algèbre symétrique de \(\underline{E}\); on notera \(\mathcal A\) la catégorie des \(\underline{O}_X\)-Modules, \(\mathcal C\) la catégorie des \(\underline{B}\)-Modules gradués avec pour morphismes les morphismes gradués, \(\mathcal C_0\) la sous-catégorie pleine de \(\mathcal C\) dont les objets sont les \(\underline{B}\)-Modules localement libres de type fini en tant que \(\underline{B}\)-Modules gradués, \(\mathcal C'\) la catégorie des \(\underline{O}_S\)-Modules. Les catégories \(\mathcal A\), \(\mathcal C\), et \(\mathcal C'\) sont abéliennes.

Pour tout \(n\), et tout complexe \(\underline{L}^\bullet\) de \(D^+(\mathcal A)\), nous noterons \(\underline{L}^\bullet(n)\) le complexe \(\underline{L}^\bullet\otimes_{\underline{O}_X}\underline{O}_X(n)\).

